\section{Introdução}

Relações comerciais fazem parte do cotidiano e envolvem, de maneira recorrente, decisões sobre preços e consumo. De modo geral, um vendedor disponibiliza seus produtos por determinados valores, enquanto os compradores decidem adquiri-los ou não com base em sua disposição a pagar. Em mercados cada vez mais competitivos, a definição de preços deixa de ser uma tarefa meramente operacional e passa a constituir uma decisão estratégica, capaz de impactar diretamente a receita obtida. Nesse contexto, surge uma questão central: como precificar um conjunto de produtos de forma a maximizar o retorno financeiro?

A resposta para essa questão não é trivial, pois a precificação envolve um equilíbrio entre atratividade e rentabilidade. Por um lado, preços mais elevados aumentam a receita obtida por unidade vendida, mas podem reduzir o conjunto de compradores dispostos a realizar a compra. Por outro, preços mais baixos tendem a ampliar a quantidade de vendas possíveis, embora reduzam o ganho individual em cada transação. Assim, do ponto de vista do vendedor, o problema consiste em encontrar uma combinação de preços que equilibre essas forças antagônicas e maximize a receita total.

Esse tipo de situação se insere em uma classe mais ampla de problemas conhecidos como \textit{Problemas de Precificação} (PP), nos quais um vendedor deve determinar os preços de seus produtos de modo a maximizar a receita total obtida com as vendas. Essa decisão é particularmente desafiadora porque o preço de cada item influencia simultaneamente o valor arrecadado por unidade e o conjunto de consumidores dispostos a adquiri-lo, evidenciando a necessidade de um equilíbrio entre competitividade e rentabilidade~\cite{smith1986}. Problemas de precificação estão presentes em diversos setores da economia e têm sido amplamente investigados na literatura, em razão de sua elevada relevância prática e aplicabilidade em contextos reais~\cite{aggarwal2004,venkatesan2005,shmuel1984,shmuel1987,smith1986}.

No contexto de problemas de precificação não paramétrica com múltiplos itens e múltiplos compradores, Rusmevichientong~\cite{rusmevichientong2006} formalizou um arcabouço geral no qual o vendedor conhece previamente as valorações dos compradores para os itens disponíveis e deve definir os preços de forma a maximizar a receita. Neste trabalho, diferentes variantes foram propostas para representar distintos comportamentos de compra. Entre elas, destacam-se o \textit{Problema da Compra Máxima} (PCmax), no qual o comprador seleciona o item viável de maior preço, e o \textit{Problema da Compra por Preferência} (PCrank), em que o comprador escolhe, dentre os itens viáveis, aquele que ocupa a posição mais alta em sua ordem de preferência.

Dentre essas variantes, este trabalho concentra-se no \textit{Problema da Compra Mínima} (PCmin). Nesta variante, o vendedor deve definir os preços de um conjunto de itens de modo a maximizar sua receita total, assumindo que cada comprador adquire, dentre os itens viáveis, aquele de menor preço. Um item é considerado viável para um comprador quando seu preço não excede a valoração atribuída a esse item por esse comprador. Caso nenhum item satisfaça essa condição, nenhuma compra é realizada. Essa hipótese de comportamento torna o PCmin particularmente relevante do ponto de vista prático, uma vez que, em diversos contextos reais, consumidores tendem a optar pelo item de menor preço dentre aqueles que consideram aceitáveis ou substitutos próximos. Embora fatores como marca, qualidade percebida e preferências individuais também possam influenciar a decisão de compra, o critério de menor preço frequentemente exerce papel determinante, especialmente em mercados com produtos de baixa diferenciação.

A Figura~\ref{fig:exemploPCmin} apresenta um exemplo ilustrativo de uma instância do PCmin. Na subfigura~(a), são representadas as relações de valoração entre os itens e os compradores por meio de um grafo bipartido, no qual cada aresta ponderada indica a valoração atribuída por um comprador a um determinado item. Na subfigura~(b), são mostrados os preços atribuídos aos itens e as escolhas induzidas pela regra do problema. Nesse caso, cada comprador adquire o item de menor preço dentre aqueles de seu interesse cujo valor não excede sua respectiva valoração. Assim, observa-se que o comprador $b_1$ seleciona o item $i_1$, o comprador $b_2$ seleciona o item $i_4$, o comprador $b_3$ seleciona o item $i_3$, enquanto o comprador $b_4$ não realiza compra, uma vez que não há item viável disponível para ele.

\begin{figure}[H]
\centering

\begin{subfigure}[t]{0.45\textwidth}
    \centering
    \conjuntoItensValorados
    \caption{Relações de valoração entre itens e compradores, representadas pelas arestas ponderadas da matriz $V$.}
\end{subfigure}
\hfill
\begin{subfigure}[t]{0.45\textwidth}
    \centering
    \conjuntoItensPrecificados
    \caption{Preços atribuídos aos itens e seleção induzida dos compradores segundo a regra do PCmin.}
\end{subfigure}

\caption{Exemplo ilustrativo de uma instância do PCmin, com as valorações entre itens e compradores e a seleção induzida pelos preços atribuídos.}
\label{fig:exemploPCmin}
\end{figure}

No exemplo ilustrado na Figura~\ref{fig:exemploPCmin}, a solução apresentada gera uma receita total de 11 unidades monetárias, correspondente às compras realizadas por $b_1$, $b_2$ e $b_3$. Entretanto, esse conjunto de preços não é necessariamente ótimo. Pequenas alterações nos preços podem modificar o conjunto de itens viáveis para cada comprador e, consequentemente, alterar tanto a quantidade de vendas quanto a receita obtida por item. Assim, elevar o preço de um item pode aumentar o ganho unitário, mas também torná-lo inviável para determinados compradores; de forma análoga, reduzir um preço pode atrair mais compradores, porém com menor retorno por venda. Esse comportamento evidencia que a definição dos preços envolve um equilíbrio não trivial entre atratividade e rentabilidade, justificando o uso de métodos de otimização para a resolução do PCmin.

No que se refere ao desenvolvimento teórico do PCmin, Rusmevichientong et al.~\cite{rusmevichientong2006} foram os primeiros a formalizar o problema no contexto de problemas de precificação não paramétrica com múltiplos itens e compradores. Além de introduzir o PCmin, os autores o inserem em uma família mais ampla de variantes de precificação, como o \textit{Problema da Compra Máxima} (PCmax) e o \textit{Problema da Compra por Preferência} (PCrank). Contudo, esse estudo possui caráter predominantemente conceitual, concentrando-se na definição do problema e em suas propriedades estruturais, sem avançar no desenvolvimento de formulações matemáticas ou métodos computacionais voltados à sua resolução prática.

Posteriormente, Briest e Krysta~\cite{briest2007} e Chalermsook et al.~\cite{chalermsook2012} aprofundaram a investigação teórica do PCmin ao estabelecer resultados de complexidade e aproximabilidade, evidenciando sua elevada dificuldade computacional. Esses trabalhos são fundamentais por consolidarem o entendimento de que o problema é NP-difícil e, portanto, desafiador sob a perspectiva algorítmica. Ainda assim, tais contribuições permanecem restritas ao campo teórico, sem propor formulações de Programação Linear Inteira Mista (PLIM), técnicas de fortalecimento ou abordagens heurísticas voltadas ao tratamento de instâncias práticas.

Mais recentemente, Catabriga~\cite{catabriga2025} representa um avanço importante ao investigar o PCmin sob a ótica da otimização exata. Em seu trabalho, são propostas diferentes formulações de Programação Linear Inteira (PLI) e Programação Linear Inteira Mista (PLIM), bem como técnicas de fortalecimento dos modelos e uma avaliação computacional em distintas classes de instâncias. Embora os resultados demonstrem o potencial dessas abordagens para a resolução do problema, também evidenciam limitações de escalabilidade em instâncias de maior porte, indicando que ainda há espaço para o fortalecimento das formulações existentes e para o desenvolvimento de estratégias complementares de solução.

Nesse contexto, o presente trabalho posiciona-se como uma continuidade dessa linha de pesquisa, com o objetivo de contribuir para o avanço do estudo do PCmin sob a perspectiva computacional. Em particular, busca-se investigar e aperfeiçoar formulações matemáticas existentes, explorar estratégias de modelagem e implementação que reduzam o tempo de execução dos métodos exatos e, adicionalmente, desenvolver e avaliar abordagens heurísticas voltadas à resolução de instâncias de grande porte em tempo hábil. Espera-se, assim, ampliar o conjunto de ferramentas disponíveis para o tratamento eficiente do PCmin, tanto no âmbito de métodos exatos quanto no de estratégias aproximativas.